% ============================================================
% CAPÍTULO 1: INTRODUÇÃO
% ============================================================

\chapter{Introdução}
\label{cap:introducao}

\section{Contextualização do Tema}
\label{sec:contextualizacao}

O desenvolvimento acelerado da inteligência artificial (IA) nas primeiras décadas do século XXI reconfigurou profundamente as relações sociais, econômicas e políticas em escala global. Sistemas algorítmicos de aprendizado de máquina, processamento de linguagem natural e reconhecimento de padrões passaram a mediar decisões que afetam diretamente a vida das pessoas --- desde a concessão de crédito e o acesso a serviços de saúde até a formação da opinião pública e a administração da justiça \cite{Floridi2023EI}. Esse fenômeno, embora portador de promessas significativas de inovação e progresso, suscita questões fundamentais sobre a proteção da pessoa humana diante de estruturas tecnológicas que operam, em larga medida, como ``caixas-pretas'' \cite{Pasquale2015BM}.

No cenário contemporâneo, a personalidade humana --- entendida como o conjunto de atributos físicos, psíquicos, morais e sociais que singularizam cada indivíduo --- enfrenta desafios inéditos. A coleta massiva de dados pessoais, a tomada de decisões automatizadas sem transparência, os vieses algorítmicos que reproduzem e amplificam desigualdades, e a crescente substituição do juízo humano por sistemas computacionais colocam em xeque categorias centrais do pensamento jurídico moderno, como a autonomia, a dignidade e a responsabilidade \cite{Zuboff2019SC, Mulholland2021IA}.

É nesse contexto que a publicação da Carta Encíclica \MH{} pelo Papa Leão XIV, em 25 de maio de 2026, representa um marco no debate ético-tecnológico contemporâneo. Dedicada à ``dignidade e à vocação da pessoa humana na era da inteligência artificial'', a Encíclica oferece uma reflexão sistemática sobre os fundamentos antropológicos e éticos que devem orientar o desenvolvimento e a regulação das tecnologias digitais. Em seu Capítulo III, particularmente, o documento aborda diretamente os desafios suscitados pela IA, propondo critérios para uma relação entre o ser humano e a máquina que seja verdadeiramente a serviço da pessoa \cite{LeaoXIV2026MH}.

O magistério pontifício insere-se, assim, em uma tradição de reflexão sobre a técnica que remonta ao Concílio Vaticano II e encontra desenvolvimento em encíclicas como a \textit{Laudato Si'} \cite{Francisco2015LS} e a \textit{Fratelli Tutti} \cite{Francisco2020FT}. No entanto, a \MH{} distingue-se por enfrentar diretamente o fenômeno da IA em suas dimensões propriamente tecnológicas, antropológicas e jurídicas, oferecendo um arsenal conceitual que pode enriquecer significativamente o debate regulatório em curso no Brasil e no mundo.

Do ponto de vista jurídico, a proteção da personalidade humana na era digital constitui um dos temas mais prementes do direito contemporâneo. O ordenamento jurídico brasileiro, alicerçado no princípio da dignidade da pessoa humana (art. 1º, III, da Constituição Federal de 1988), dispõe de instrumentos como os direitos da personalidade (arts. 11 a 21 do Código Civil), a Lei Geral de Proteção de Dados Pessoais (Lei nº 13.709/2018) e o Marco Civil da Internet (Lei nº 12.965/2014). Contudo, a velocidade da inovação tecnológica desafia a capacidade do direito de acompanhar e regular adequadamente as novas realidades, gerando lacunas normativas e incertezas interpretativas que demandam fundamentos sólidos para a construção de respostas jurídicas coerentes \cite{MendesDoneda2021PD, Frazao2020RD}.

A presente monografia insere-se nessa intersecção entre direito, tecnologia e ética, propondo-se a investigar as contribuições da \MH{} para a proteção jurídica da personalidade humana na era da IA.

\section{Problema de Pesquisa}
\label{sec:problema}

Considerando o cenário descrito, a pesquisa parte do seguinte problema central:

\begin{citacao}
\textbf{Como os fundamentos antropológicos e éticos propostos pela Encíclica \MH{} podem subsidiar a construção de um regime jurídico de proteção da personalidade humana diante dos desafios suscitados pela inteligência artificial no direito contemporâneo brasileiro?}
\end{citacao}

O problema desdobra-se nas seguintes questões orientadoras:

\begin{enumerate}
    \item Qual é a concepção de pessoa humana subjacente ao magistério da \MH{} e como ela se articula com a tradição do pensamento personalista?
    \item Quais são os principais desafios que a inteligência artificial impõe à proteção da personalidade humana?
    \item Como o ordenamento jurídico brasileiro atual protege a personalidade humana diante das tecnologias de IA?
    \item De que forma os princípios éticos propostos pela \MH{} podem orientar a regulação da IA e o aperfeiçoamento da proteção jurídica da personalidade?
\end{enumerate}

\section{Hipóteses}
\label{sec:hipoteses}

A pesquisa se orienta pelas seguintes hipóteses:

\begin{description}
    \item[H1:] \textbf{Hipótese antropológica.} A antropologia personalista da \MH{}, fundada na dignidade ontológica, relacional e transcendental da pessoa humana, oferece um referencial teórico mais robusto para a proteção jurídica da personalidade na era digital do que as abordagens exclusivamente instrumentais ou liberais-individualistas.

    \item[H2:] \textbf{Hipótese crítica.} O paradigma tecnocrático criticado pela \MH{} --- que reduz o ser humano a dados e a decisões à eficiência algorítmica --- constitui uma ameaça real e imediata à integridade da personalidade humana, que o direito contemporâneo ainda não enfrenta adequadamente.

    \item[H3:] \textbf{Hipótese normativa.} Os princípios éticos da subsidiariedade, solidariedade, transparência e responsabilidade, conforme reelaborados pela \MH{}, podem ser traduzidos em critérios normativos operacionais para a regulação da IA e para a interpretação dos institutos jurídicos existentes.

    \item[H4:] \textbf{Hipótese propositiva.} A incorporação dos fundamentos da \MH{} ao debate jurídico brasileiro sobre regulação da IA pode contribuir para superar as limitações da abordagem exclusivamente técnica ou economicista que predomina nos atuais projetos legislativos.
\end{description}

\section{Objetivos}
\label{sec:objetivos}

\subsection{Objetivo Geral}
\label{subsec:objetivo-geral}

O objetivo geral da pesquisa é \textbf{analisar as contribuições da Encíclica \MH{} para a construção de um regime jurídico de proteção da personalidade humana diante dos desafios suscitados pela inteligência artificial, à luz do direito contemporâneo brasileiro}.

\subsection{Objetivos Específicos}
\label{subsec:objetivos-especificos}

\begin{enumerate}
    \item[OE1.] Examinar os fundamentos antropológicos e éticos da \MH{}, identificando sua concepção de pessoa humana e sua relação com a tradição do personalismo cristão e da doutrina social da Igreja.
    \item[OE2.] Identificar e analisar os principais desafios que a inteligência artificial impõe à proteção da personalidade humana, com base na literatura especializada e no diagnóstico oferecido pela Encíclica.
    \item[OE3.] Mapear o regime jurídico brasileiro de proteção da personalidade, com ênfase nos direitos da personalidade, na LGPD e no Marco Civil da Internet, avaliando sua adequação diante dos desafios da IA.
    \item[OE4.] Extrair da \MH{} critérios normativos que possam orientar a regulação da IA e a interpretação dos institutos jurídicos de proteção da personalidade.
    \item[OE5.] Demonstrar como os fundamentos da \MH{} podem contribuir para o aperfeiçoamento do debate regulatório brasileiro sobre IA, em particular o PL 2.338/2023.
\end{enumerate}

\section{Justificativa}
\label{sec:justificativa}

A justificativa para a presente pesquisa assenta-se em três ordens de razões: acadêmica, social e jurídica.

No plano \textbf{acadêmico}, verifica-se uma lacuna na literatura jurídica brasileira no que tange ao diálogo entre o magistério pontifício recente e o direito da inteligência artificial. Embora existam estudos sobre ética da IA, direitos da personalidade e proteção de dados, são escassas as pesquisas que se debruçam sistematicamente sobre as contribuições do pensamento social cristão para a regulação da tecnologia. A publicação da \MH{}, por sua atualidade e abrangência, oferece uma oportunidade ímpar para preencher essa lacuna.

No plano \textbf{social}, a pesquisa responde a uma demanda concreta da sociedade brasileira por fundamentos éticos sólidos que orientem o desenvolvimento tecnológico. Em um contexto de crescente digitalização das relações sociais e de concentração do poder tecnológico em poucas corporações, a reflexão sobre a proteção da pessoa humana diante da IA transcende o interesse acadêmico e se configura como uma necessidade democrática.

No plano \textbf{jurídico}, a pesquisa se justifica pela urgência de se construir um regime de regulação da IA que seja coerente com os fundamentos constitucionais brasileiros, em especial o princípio da dignidade da pessoa humana. A tramitação do PL 2.338/2023 no Congresso Nacional e a necessidade de adequação do ordenamento jurídico brasileiro aos parâmetros internacionais (como o AI Act europeu) tornam premente a reflexão sobre os critérios ético-jurídicos que devem informar essa regulação.

\section{Metodologia}
\label{sec:metodologia}

A presente pesquisa adota metodologia de natureza \textbf{jurídico-teórica}, com abordagem \textbf{qualitativa} e método de procedimento \textbf{hermenêutico-crítico}.

Quanto à \textbf{natureza}, a pesquisa é jurídico-teórica porque se dedica à análise de conceitos, princípios e fundamentos do direito, sem pretensão de construir generalizações empíricas, mas buscando compreender o sentido e o alcance das normas e dos discursos jurídicos à luz de um referencial teórico determinado \cite{Comparato2016DH}.

Quanto à \textbf{abordagem}, a pesquisa é qualitativa, voltada para a interpretação aprofundada de textos normativos, documentos e doutrina, privilegiando a compreensão dos significados e das relações entre os conceitos em detrimento da quantificação de dados.

Quanto ao \textbf{método}, adota-se a abordagem hermenêutico-crítica, que combina a compreensão interpretativa dos textos (hermenêutica) com a análise das estruturas de poder e das condições sociais que subjazem aos discursos (crítica). Esse método é particularmente adequado para investigar textos de natureza ético-jurídica, como uma encíclica pontifícia e as normas de proteção da personalidade, pois permite simultaneamente compreender seu sentido interno e avaliar sua capacidade de responder aos desafios do contexto social \cite{Ricoeur1990SM}.

As \textbf{técnicas de pesquisa} empregadas são:

\begin{enumerate}
    \item \textbf{Pesquisa bibliográfica sistematizada}, com levantamento de obras de referência nos campos do direito constitucional, direitos da personalidade, proteção de dados, ética da inteligência artificial e doutrina social da Igreja.
    \item \textbf{Análise documental}, compreendendo o exame integral do texto da Encíclica \MH{}, da legislação brasileira pertinente (CF/88, CC/02, LGPD, Marco Civil da Internet), do PL 2.338/2023 e do AI Act europeu.
    \item \textbf{Levantamento jurisprudencial}, com ênfase na jurisprudência do STJ e do STF sobre direitos da personalidade e proteção de dados.
\end{enumerate}

O \textbf{plano de análise} segue a estruturação em capítulos da monografia, articulando progressivamente os fundamentos teóricos (capítulos 2 e 3), o diagnóstico jurídico-positivo (capítulo 4) e a proposição de critérios normativos (capítulo 5), culminando na síntese conclusiva (capítulo 6).

\section{Delimitação do Estudo}
\label{sec:delimitacao}

A pesquisa circunscreve-se ao exame da proteção jurídica da personalidade humana no direito brasileiro, com ênfase nos aspectos constitucionais e infraconstitucionais (Código Civil, LGPD, Marco Civil da Internet), à luz dos fundamentos antropológicos e éticos propostos pela Encíclica \MH{}. Não se pretende, portanto, oferecer uma teoria geral da regulação da IA, tampouco esgotar as múltiplas implicações tecnológicas, econômicas ou políticas do fenômeno. A análise concentra-se na dimensão propriamente jurídico-antropológica do problema, investigando como uma determinada concepção de pessoa humana pode informar e qualificar a proteção jurídica diante dos desafios algorítmicos.

No que tange ao direito comparado, a pesquisa limita-se a referências pontuais ao AI Act europeu e ao GDPR, na medida em que estes influenciam diretamente o debate legislativo brasileiro, sem pretensão de realizar um estudo comparativo exaustivo.

\section{Estrutura da Monografia}
\label{sec:estrutura}

A presente monografia estrutura-se em seis capítulos, além das referências e anexos.

O \textbf{Capítulo 1} (Introdução) apresenta a contextualização do tema, o problema de pesquisa, as hipóteses, os objetivos, a justificativa e a metodologia adotada.

O \textbf{Capítulo 2} (Fundamentos Antropológicos e Éticos da \MH{}) examina a concepção de pessoa humana subjacente ao magistério pontifício, situando a Encíclica na tradição do personalismo cristão e analisando os princípios éticos que orientam sua reflexão sobre a técnica.

O \textbf{Capítulo 3} (IA e os Desafios à Personalidade Humana) identifica e analisa as principais ameaças que as tecnologias de IA representam para a integridade da personalidade humana, com base no diagnóstico oferecido pela \MH{} e na literatura especializada.

O \textbf{Capítulo 4} (Proteção Jurídica da Personalidade no Direito Brasileiro) mapeia o regime jurídico de proteção da personalidade, avaliando sua adequação diante dos desafios da IA.

O \textbf{Capítulo 5} (Contribuições da \MH{} para a Regulação da IA) propõe a tradução dos fundamentos éticos da Encíclica em critérios normativos operacionais para a regulação da IA e para o aperfeiçoamento da proteção jurídica da personalidade.

O \textbf{Capítulo 6} (Conclusão) retoma as hipóteses, sintetiza as contribuições da pesquisa, aponta suas limitações e sugere agendas para investigações futuras.
